\newpage
\section*{Pr\'actica 3}
\subsection*{1: Popurri de ejercicios sobre probabilidades condicionales y el teorema de Bayes.}
\begin{enumerate}
    \item En un bolso hay una pelota. Ignoro si es blanca o negra. Meto una segunda pelota dentro del bolso, que es de color blanco. Sacudo el bolso, meto la mano y extraigo una pelota blanca. Cu\'al es la probabilidad de que la pelota que queda sea blanca?
    \item Ana tiene dos hermanas: Beatriz y Camila. Ana es m\'as grande que Beatriz. Que probaiblidad hay de que Ana sea m\'as grande que Camila?
    \item En un caj\'on de mi casa hay 4 facturas de gas y una de luz para pagar. Mi mujer abre el caj\'on, saca una factura al azar y la paga sin decirme cu\'a es. En ese momento, llega una nueva factura de luz, que va a para al caj\'on de las facturas impagas. Abro el caj\'on y tomo una factura al azar. Es de gas. Qu\'e probabilidad hay de que mi mujer haya pagado una factura de gas?
\end{enumerate}

\subsection*{2: La pequeña ameba.}
Una ameba, Bobo, vive en un pantano. Cada minuto a cada ameba le sucede de forma independiente una de las siguientes con probabilidades iguales:
\begin{enumerate}
    \item La ameba muere;
    \item La ameba se mantiene igual o
    \item La ameba se duplica.
\end{enumerate}
Calcule la probabilidad de que el linaje de Bobo se extinga.
\subsection*{3: Un caso parad\'ogico: el valor de la informaci\'on}
Realice los siguientes calculos frente a la informaci\'on brindada
\begin{enumerate}
    \item {\it ``Tengo dos hijos, la primera es mujer''}. Cu\'al es la probabilidad de que la otra tambi\'en lo sea?
    \item {\it ``Tengo dos hijos, una es mujer''}. Cu\'al es la probabilidad de que la otra tambi\'en?
    \item {\it ``Tengo dos hijos, una es mujer y naci\'o un martes''}.  Cu\'al es la probabilidad de que la otra tambi\'en?
\end{enumerate}
\subsection*{4: Un ejemplo hist\'orico: el valor del contexto}
En un juego televisivo hay 3 puertas cerradas, y un premio se esconde detr\'as de una de las puertas. El participante elige una puerta al azar. El conductor entonces abre alguna de las dos puertas no elegidas por el participante, cuidando de dejar cerrada aquella que contiene el premio. El premio ahora puede estar en la puerta elegida inicialmente por el participante, o en la otra puerta, que el conductor dej\'o cerrada. Cu\'al es la probabilidad de que el premio se encuentre en cada una de ellas?

En una de las realizaciones del juego, cuando el conductor est\'a por abrir una de las puertas, es interrumpido por un terremoto. Este no es muy fuerte pero abre una de las puertas accidentalmente. Al abrirse muestra que detr\'as de ella no est\'a el premio. A lo que el entrevistador dice {\it ``Da lo mismo, es como si la hubiera abierto yo''} y decide seguir con el programa. Cu\'al es la probabilidad de que el premio se encuentre en cada una de las puertas?


\subsection*{5: Inferiendo la probabilidad de cara de una moneda}
Te encontraste una moneda en la calle. Tu transtorno obsesivo compulsivo te obliga a determinar la probabilidad de que salga cara al arrojarla. Para eso arrojas la moneda $N$ veces. C\'omo usas la informaci\'on del resultado para realizar la inferencia? Utilice todas las herramientas que consideres necesarias. Considere su grado de incerteza en su estimaci\'on. Luego, calcule la probabilidad de obtener cara en el siguiente tiro de la moneda.